\documentclass[12pt]{article}
\usepackage[utf8]{inputenc}
\usepackage[T1]{fontenc}
\usepackage[margin=1in]{geometry}
\usepackage{setspace}
\usepackage{xcolor}
\usepackage[colorlinks=true, urlcolor=blue!60!black, linkcolor=blue!60!black]{hyperref}
\setlength{\parindent}{2em}
\setlength{\parskip}{1.25em}
\renewcommand{\baselinestretch}{1.0}
\title{\textit{Flast v. Cohen}, 392 U.S. 83 (1968)}
\author{Case Brief}
\date{}
\begin{document}
\maketitle

\section{Case Metadata}
\begin{itemize}
  \item \textbf{Case Name:} \textit{Flast v. Cohen}
  \item \textbf{Citation:} 392 U.S. 83 (1968) (No. 416); \href{https://supreme.justia.com/cases/federal/us/392/83/}{Justia Opinion}; \href{https://www.law.cornell.edu/constitution-conan/article-3/section-2/clause-1/taxpayer-standing}{Cornell LII Constitution Annotated}
  \item \textbf{Dates:} Argued March 12, 1968; Decided June 10, 1968
  \item \textbf{Lower Court:} Three-judge U.S. District Court for the Southern District of New York, 271 F. Supp. 1 (S.D.N.Y. 1967) (dismissed for lack of standing under \textit{Frothingham v. Mellon}; Judge Frankel dissenting; direct appeal under 28 U.S.C. \S\ 1253)
  \item \textbf{Vote:} 8--1 (Reversed)
  \item \textbf{Author:} Chief Justice Warren (joined by Justices Black, Douglas, Brennan, Stewart, White, Fortas, and Marshall)
  \item \textbf{Concurrences:} Justice Douglas; Justice Stewart; Justice Fortas
  \item \textbf{Dissent:} Justice Harlan
  \item \textbf{Constitutional Provisions:} Article III, \S\ 2 (Case-or-Controversy Standing); Article I, \S\ 8, cl. 1 (Taxing and Spending Clause); First Amendment Establishment Clause
\end{itemize}

\section{Facts and Procedural History}
Seven federal income taxpayers, led by Florence Flast, sued Secretary of Health, Education, and Welfare Wilbur Cohen and other federal officials in the Southern District of New York to enjoin federal expenditures under Titles I and II of the Elementary and Secondary Education Act of 1965 (ESEA), which appropriated nearly \$1 billion in federal grants to state and local educational agencies (\href{https://supreme.justia.com/cases/federal/us/392/83/}{Justia Opinion}). Title I funded special educational services for low-income children, including students enrolled in private schools; Title II funded school library resources, textbooks, and instructional materials for public and private elementary and secondary schools.

Plaintiffs rested their standing solely on their status as federal income taxpayers. They alleged that federal funds were being disbursed to finance instruction in reading, arithmetic, and other subjects and to purchase textbooks and library materials for religious and sectarian schools, violating the Establishment and Free Exercise Clauses of the First Amendment. They sought a declaration that the expenditures were unauthorized by the ESEA or, alternatively, unconstitutional, along with an injunction.

Relying on \textit{Frothingham v. Mellon}, 262 U.S. 447 (1923)---which had barred federal taxpayer standing for 45 years on the ground that an individual taxpayer's interest in the Federal Treasury is ``comparatively minute and indeterminable''---the Government moved to dismiss. Judge Frankel convened a three-judge District Court (267 F. Supp. 351 (S.D.N.Y. 1967)), which dismissed the complaint on \textit{Frothingham}'s authority over Judge Frankel's dissent (271 F. Supp. 1 (S.D.N.Y. 1967)). Plaintiffs took a direct appeal to the Supreme Court.

\section{Issue(s) Presented}
Whether Article III or \textit{Frothingham v. Mellon} bars federal taxpayers, suing solely in their capacity as taxpayers, from establishing standing to challenge federal expenditures under the Elementary and Secondary Education Act of 1965 as violating the Establishment Clause of the First Amendment.

\section{Holding}
No. In an 8--1 decision, the Supreme Court reversed and carved out a narrow ``double-nexus'' exception to \textit{Frothingham} (\href{https://supreme.justia.com/cases/federal/us/392/83/}{Justia Opinion}). Federal taxpayers have Article III standing to challenge a federal statute when they establish (1) a logical link between their taxpayer status and an exercise of congressional power under Article I, \S\ 8's Taxing and Spending Clause, and (2) a nexus between their taxpayer status and a specific constitutional limitation imposed upon the taxing and spending power---such as the Establishment Clause.

\section{Reasoning}
After confirming direct appellate jurisdiction because a three-judge court was properly convened under 28 U.S.C. \S\ 2282, Chief Justice Warren reexamined \textit{Frothingham} and Article III justiciability (\href{https://supreme.justia.com/cases/federal/us/392/83/}{Justia Opinion}).

\subsection{Reinterpreting Article III Standing and \textit{Frothingham}}
Warren noted that commentators long debated whether \textit{Frothingham} announced an absolute Article III bar or a prudential rule of self-restraint. Rejecting the Government's claim that separation of powers categorically forecloses taxpayer suits, Warren drew a sharp line between standing and other justiciability doctrines: standing focuses solely on whether the litigant is a ``proper party'' with ``such a personal stake in the outcome of the controversy as to assure that concrete adverseness'' required by \textit{Baker v. Carr}, 369 U.S. 186, 204 (1962). Separation-of-powers concerns arise from the substantive issues presented (such as political questions), not from the identity of the plaintiff. Consequently, Article III imposes no absolute bar to federal taxpayer suits.

\subsection{The Double-Nexus Test}
To distinguish proper taxpayer suits from impermissible ``generalized grievances,'' Warren announced a two-part nexus test:
\begin{itemize}
  \item \textbf{First Nexus (Type of Enactment):} The taxpayer must challenge an exercise of congressional power under the Taxing and Spending Clause of Article I, \S\ 8---not merely an ``incidental expenditure of tax funds in the administration of an essentially regulatory statute.''
  \item \textbf{Second Nexus (Nature of Constitutional Limit):} The taxpayer must allege that the challenged enactment exceeds a \textit{specific constitutional limitation} imposed upon Congress's taxing and spending power, rather than merely asserting that Congress acted generally beyond its delegated Article I, \S\ 8 powers.
\end{itemize}

Applying this framework, the Court held that Flast satisfied both prongs. First, Titles I and II of the ESEA represented a massive spending program enacted under Article I, \S\ 8. Second, invoking James Madison's \textit{Memorial and Remonstrance Against Religious Assessments} (warning against forcing citizens to contribute ``three pence'' to a religious establishment), Warren concluded that the Establishment Clause was designed specifically to restrict Congress's taxing and spending power from aiding religion. By contrast, \textit{Frothingham} survived intact: Mrs. Frothingham satisfied the first nexus by challenging a federal spending statute (the Maternity Act of 1921), but failed the second because she invoked only the Tenth Amendment and Fifth Amendment Due Process Clause---neither of which specifically restricts federal spending.

\section{Concurring and Dissenting Opinions}
\begin{itemize}
  \item \textbf{Justice Douglas (Concurring):} Joined the opinion but predicted that the majority's double-nexus test would prove unsustainable for the very reasons Justice Harlan gave (\href{https://supreme.justia.com/cases/federal/us/392/83/}{Justia Opinion}). Urging the Court to overrule \textit{Frothingham} outright, Douglas argued that taxpayers should serve as vigilant ``private attorneys general'' to enforce all constitutional guarantees against federal overreach.
  \item \textbf{Justice Stewart (Concurring):} Emphasized that the Court held \textit{only} that taxpayers have standing to enforce the Establishment Clause's explicit personal right not to be taxed to support religion, while preserving \textit{Frothingham}'s bar against generalized grievances (\href{https://supreme.justia.com/cases/federal/us/392/83/}{Justia Opinion}).
  \item \textbf{Justice Fortas (Concurring):} Agreed that \textit{Flast} must be strictly confined to the Establishment Clause and should never become a ``launching pad'' for taxpayer attacks on other federal spending programs (\href{https://supreme.justia.com/cases/federal/us/392/83/}{Justia Opinion}).
  \item \textbf{Justice Harlan (Dissenting):} Delivered a classic critique of the double-nexus test (\href{https://supreme.justia.com/cases/federal/us/392/83/}{Justia Opinion}). Harlan argued that taxpayers challenging general appropriations are ``non-Hohfeldian'' plaintiffs asserting no personal or proprietary stake, because tax revenues merge into general funds and winning does not reduce a plaintiff's tax bill. Neither prong of Warren's test actually measures a plaintiff's personal stake: a citizen's interest does not depend on whether an expenditure is ``incidental'' to regulation or whether a clause is labeled a ``specific limitation.'' While Article III does not strictly forbid Congress from authorizing ``public actions,'' Harlan argued that courts should not invent unprincipled taxpayer standing without congressional authorization lest the judiciary become a ``Council of Revision.''
\end{itemize}

\section{Key Quotations}
\begin{quote}
``First, the taxpayer must establish a logical link between that status and the type of legislative enactment attacked. Thus, a taxpayer will be a proper party to allege the unconstitutionality only of exercises of congressional power under the taxing and spending clause of Art. I, \S\ 8, of the Constitution. ... Secondly, the taxpayer must establish a nexus between that status and the precise nature of the constitutional infringement alleged. Under this requirement, the taxpayer must show that the challenged enactment exceeds specific constitutional limitations imposed upon the exercise of the congressional taxing and spending power, and not simply that the enactment is generally beyond the powers delegated to Congress by Art. I, \S\ 8.'' --- Chief Justice Warren, \href{https://supreme.justia.com/cases/federal/us/392/83/}{392 U.S. at 102--103}
\end{quote}

\begin{quote}
``The taxpayer's allegation in such cases would be that his tax money is being extracted and spent in violation of specific constitutional protections against such abuses of legislative power.'' --- Chief Justice Warren, \href{https://supreme.justia.com/cases/federal/us/392/83/}{392 U.S. at 106}
\end{quote}

\begin{quote}
``The status of taxpayer should not be accepted as a launching pad for an attack upon any target other than legislation affecting the Establishment Clause.'' --- Justice Fortas, concurring, \href{https://supreme.justia.com/cases/federal/us/392/83/}{392 U.S. at 116}
\end{quote}

\begin{quote}
``To describe those rights and interests as personal, and to intimate that they are in some unspecified fashion to be differentiated from those of the general public, reduces constitutional standing to a word game played by secret rules.'' --- Justice Harlan, dissenting, \href{https://supreme.justia.com/cases/federal/us/392/83/}{392 U.S. at 129}
\end{quote}

\section{Doctrinal Significance and Modern Legacy (Through 2026)}
Although never formally overruled, \textit{Flast v. Cohen} has been tightly cabined into a near-singular anomaly in modern Article III standing doctrine (\href{https://www.law.cornell.edu/constitution-conan/article-3/section-2/clause-1/taxpayer-standing}{Cornell LII Constitution Annotated}). Outside one narrow as-applied grant challenge in \href{https://supreme.justia.com/cases/federal/us/487/589/}{\textit{Bowen v. Kendrick}}, 487 U.S. 589 (1988) (allowing taxpayer standing where Congress expressly authorized grants to religious organizations under the Adolescent Family Life Act), the Supreme Court has rejected every effort to extend \textit{Flast}:
\begin{itemize}
  \item \textbf{Rejection Beyond the Establishment Clause (1974):} In \href{https://supreme.justia.com/cases/federal/us/418/166/}{\textit{United States v. Richardson}}, 418 U.S. 166 (1974) (Statement and Account Clause challenge to secret CIA budgets), and \href{https://supreme.justia.com/cases/federal/us/418/208/}{\textit{Schlesinger v. Reservists Committee to Stop the War}}, 418 U.S. 208 (1974) (Incompatibility Clause challenge), the Court refused to recognize any other ``specific constitutional limitation'' under \textit{Flast}.
  \item \textbf{Property Clause and Executive Transfers (\textit{Valley Forge}, 1982):} In \href{https://supreme.justia.com/cases/federal/us/454/464/}{\textit{Valley Forge Christian College v. Americans United for Separation of Church and State, Inc.}}, 454 U.S. 464 (1982), the Court denied taxpayer standing to challenge a federal agency's free transfer of surplus government land to a religious college because the transfer was an \textit{executive} action taken under Congress's Article IV Property Clause rather than an Article I, \S\ 8 spending enactment.
  \item \textbf{Discretionary Executive Spending (\textit{Hein}, 2007):} In \href{https://supreme.justia.com/cases/federal/us/551/587/}{\textit{Hein v. Freedom From Religion Foundation, Inc.}}, 551 U.S. 587 (2007), a 5--4 Court denied taxpayer standing to challenge White House Faith-Based Initiatives conferences funded out of general executive-branch appropriations rather than an express congressional mandate. Justice Alito's plurality declared that \textit{Flast} is ``neither extended nor overruled,'' while Justices Scalia and Thomas urged overruling \textit{Flast} as incoherent.
  \item \textbf{Tax Credits vs. Expenditures (\textit{Winn}, 2011):} In \href{https://supreme.justia.com/cases/federal/us/563/125/}{\textit{Arizona Christian School Tuition Organization v. Winn}}, 563 U.S. 125 (2011), the Court held 5--4 that taxpayers lacked standing under \textit{Flast} to challenge state tax credits for contributions to religious school scholarship organizations because a tax credit allows private citizens to keep their own money rather than ``extracting and spending'' dissenting taxpayers' funds.
  \item \textbf{Status Through 2026:} Through 2026, \textit{Flast} survives only as a razor-thin relic confined strictly to direct Establishment Clause challenges against express congressional appropriations under Article I, \S\ 8. The modern Roberts Court's strict Article III injury-in-fact jurisprudence---reaffirmed across standing decisions through 2026, including \href{https://supreme.justia.com/cases/federal/us/609/26a203/}{\textit{National Park Service v. National Trust for Historic Preservation}}, 609 U.S. \_\_\_ (2026), rejecting psychological distress as a concrete injury---leaves \textit{Flast} completely isolated from general standing law.
\end{itemize}

\section{Pedagogical Critiques and Exam Hypotheticals}
\subsection{Key Critiques}
\begin{itemize}
  \item \textbf{Formalism and Arbitrariness:} Critics from both sides---from Justice Harlan to Justice Scalia in \textit{Hein} and Justice Kagan dissenting in \textit{Winn}---have noted the economic and conceptual arbitrariness of granting standing when Congress spends \$1 million directly on religious schools (\textit{Flast}), yet denying standing when the President spends \$1 million from general appropriations (\textit{Hein}), transfers \$1 million in surplus land (\textit{Valley Forge}), or grants \$1 million in tax credits (\textit{Winn}).
  \item \textbf{Tension with \textit{Lujan} Separation of Powers:} Chief Justice Warren's claim that standing does not implicate separation of powers has been repudiated by modern Article III doctrine (\textit{Lujan v. Defenders of Wildlife}, 504 U.S. 555 (1992)), which treats concrete, particularized injury as an essential separation-of-powers boundary.
\end{itemize}

\subsection{Exam Hypotheticals and ``Opposite Side'' Contrast}
\textbf{Hypothetical 1 (Express Congressional Appropriation to Houses of Worship):} Congress passes the ``Community Heritage Act'' under Article I, \S\ 8, expressly appropriating \$50 million to the Department of Housing and Urban Development with a statutory directive to distribute direct repair grants to churches and synagogues for sanctuary renovations. A federal income taxpayer with no other connection to the program sues in federal court under the Establishment Clause. \textit{Analysis:} The taxpayer \textit{has standing} under \textit{Flast} and \href{https://supreme.justia.com/cases/federal/us/487/589/}{\textit{Bowen v. Kendrick}}. Both nexuses are met: (1) the challenge targets an express congressional spending enactment under Article I, \S\ 8 (not general executive overhead or a regulatory statute), and (2) the challenge invokes the Establishment Clause, the sole recognized specific constitutional limitation on taxing and spending.

\textbf{Hypothetical 2 (``Opposite Side'' Contrast --- Outside \textit{Flast}'s Razor-Thin Window):} Suppose instead that (a) a federal taxpayer challenges a \$50 million congressional appropriation funding foreign military aid as violating Article I, \S\ 8's Declare War Clause or the Tenth Amendment, \textit{or} (b) a federal taxpayer brings an Establishment Clause challenge against either a new federal tax credit for donations to religious charities or a cabinet secretary's use of general agency operating funds to host sectarian prayer rallies. \textit{Contrast:} In every one of these scenarios, the suit must be dismissed for lack of Article III standing. Scenario (a) fails \textit{Flast}'s second nexus under \textit{Frothingham}, \href{https://supreme.justia.com/cases/federal/us/418/166/}{\textit{Richardson}}, and \href{https://supreme.justia.com/cases/federal/us/418/208/}{\textit{Schlesinger}} because no clause other than the Establishment Clause qualifies as a specific spending limitation. Scenario (b) fails \textit{Flast}'s first nexus under \href{https://supreme.justia.com/cases/federal/us/563/125/}{\textit{Winn}} (tax credits do not ``extract and spend'' taxpayer funds) and \href{https://supreme.justia.com/cases/federal/us/551/587/}{\textit{Hein}} (discretionary executive spending from general appropriations is not an express congressional spending mandate).

\end{document}
